Let
\[
 \mathcal E=\mathcal V\cup\mathcal R,
 \qquad
 L=\left|\mathcal X_{\mathcal V}\times\mathcal X_{\mathcal R}\right|,
 \qquad
 D=(L-1)+|\mathcal X_A|(|\mathcal X_Y|-1),
\]
and put
\[
 K^{\sharp}=D+1=L+|\mathcal X_A|(|\mathcal X_Y|-1).
\]
The finite-alphabet assumption makes all four quantities finite.  We prove the stronger intermediate assertion that every factor needs at most \(K^{\sharp}\) support points.

We first derive the finite-dimensional support fact used below.  Let \((\Omega,\mu)\) be a probability space, let \(z:\Omega\to\mathbb R^m\) be bounded and measurable, and suppose its range lies in an affine subspace of dimension at most \(D\).  Set \(b=\int z\,d\mu\).  We claim that there are \(k\leq D+1\), points \(u_1,\ldots,u_k\in\Omega\), and weights \(\alpha_1,\ldots,\alpha_k\geq0\) such that
\[
 \sum_{i=1}^k\alpha_i=1,
 \qquad
 b=\sum_{i=1}^k\alpha_i z(u_i).
 \tag{1}
\]
Here is a proof, included so that no undeclared convex-geometric result is used.  Write \(S=z(\Omega)\), \(C=\operatorname{conv}(S)\), and \(K=\overline C\).  Partition a bounded box containing \(S\) into finitely many Borel sets of diameter at most \(1/n\), and, from every set having nonempty inverse image, choose one value of \(z\) in that set.  Replacing \(z\) by the resulting finite-valued measurable map gives \(z_n\) with \(\|z_n-z\|\leq1/n\) and \(\int z_n\,d\mu\in C\).  Hence \(b\in K\).

We next show \(b\in C\) by induction on \(d=\dim\operatorname{aff}(K)\).  The assertion is immediate for \(d=0\).  If \(b\) is in the relative interior of \(K\), a sufficiently small \(d\)-simplex centered at \(b\) lies in \(K\).  Approximate its vertices by points of \(C\).  For sufficiently accurate approximations, \(b\) remains in the interior of the approximating simplex, so \(b\in C\).  If \(b\) is on the relative boundary of \(K\), the elementary supporting-hyperplane argument supplies a nonconstant affine functional \(\ell\) on \(\operatorname{aff}(K)\) such that
\[
 \ell(x)\geq\ell(b)\quad(x\in K).
\]
For completeness, this functional is obtained by taking points outside \(K\) converging to \(b\), projecting them onto the closed convex set \(K\), normalizing the projection residuals, and passing to a convergent subsequence in the finite-dimensional unit sphere.  Linearity of integration now gives
\[
 0=\ell(b)-\ell(b)
 =\int\{\ell(z(u))-\ell(b)\}\,d\mu(u).
\]
The integrand is nonnegative, so it vanishes almost surely.  After restricting to that full-measure set, the range of \(z\) lies in the proper affine face \(K\cap\{\ell=\ell(b)\}\).  The induction hypothesis on its strictly smaller affine dimension yields \(b\in C\).

Thus \(b\) is some finite convex combination of points of \(S\).  If it uses more than \(D+1\) points with positive weights, their augmented vectors \((1,z(u_i))\) are linearly dependent because the affine dimension is at most \(D\).  Hence there are numbers \(c_i\), not all zero, satisfying
\[
 \sum_i c_i=0,
 \qquad
 \sum_i c_i z(u_i)=0.
\]
After changing the common sign if necessary, some \(c_i>0\).  With
\[
 t=\min_{i:c_i>0}\frac{\alpha_i}{c_i},
 \qquad
 \alpha_i'=\alpha_i-tc_i,
\]
all \(\alpha_i'\) are nonnegative, at least one is zero, their sum is one, and their barycenter is still \(b\).  Repetition proves (1).

We apply (1) factor by factor.  By the semi-Markovian assumption, the exogenous factors indexed by
\[
 \mathcal J=\mathcal E^{\leftrightarrow}(\mathcal G)\sqcup\mathcal V\sqcup\mathcal R
\]
have a product law \(\bigotimes_{h\in\mathcal J}\mu_h\).  The graph is finite, so \(\mathcal J\) is finite.  By acyclicity, the structural equations have a unique recursive deterministic evaluation \(\Phi(u)\in\mathcal X_{\mathcal V}\times\mathcal X_{\mathcal R}\) for each joint factor state \(u\).  For \(a\in\mathcal X_A\), let \(\Phi^a(u)\) be the corresponding evaluation after replacing only the equation for \(A\) by the constant \(a\).  This intervention retains the structural equations for every missingness indicator, as required by the natural interventional law in the causal-query definition.

Fix a factor \(h\in\mathcal J\), temporarily keeping the structural equations and all other factor laws fixed, and write \(\mu_{-h}=\bigotimes_{j\neq h}\mu_j\).  Define the bounded measurable vector
\[
 \begin{aligned}
 z_h(u_h)=\Bigg(&
 \left(
   \int \mathbf 1\{\Phi(u_h,u_{-h})=(v,r)\}\,d\mu_{-h}(u_{-h})
 \right)_{(v,r)},\\
 &\left(
   \int \mathbf 1\{\Phi_Y^a(u_h,u_{-h})=y\}\,d\mu_{-h}(u_{-h})
 \right)_{(a,y)}
 \Bigg).
 \end{aligned}
 \tag{2}
\]
The first block in (2) is a probability vector on the \(L\) full cells.  For each fixed \(a\), the corresponding outcome block is a probability vector on \(\mathcal X_Y\).  Consequently the range of \(z_h\) obeys the independent affine equations
\[
 \sum_{(v,r)}z_{h,v,r}=1,
 \qquad
 \sum_{y\in\mathcal X_Y}z_{h,a,y}=1
 \quad(a\in\mathcal X_A),
\]
and therefore lies in an affine subspace of dimension at most
\[
 \{L+|\mathcal X_A||\mathcal X_Y|\}
 -\{1+|\mathcal X_A|\}
 =D.
 \tag{3}
\]

Mutual independence and Tonelli's theorem for the nonnegative indicators in (2) give
\[
 \int z_h(u_h)\,d\mu_h(u_h)
 =\left(
   \bigl(P(V=v,R=r)\bigr)_{(v,r)},
   \bigl(P(Y=y\mid\operatorname{do}(A=a))\bigr)_{(a,y)}
 \right).
 \tag{4}
\]
The second block equals \((Q_a^{\mathcal M}(y))_{(a,y)}\) by the causal-query definition.  Applying (1) with the dimension bound (3) replaces \(\mu_h\) by an atomic probability measure \(\mu_h'\), supported on at most \(K^{\sharp}\) actual values of the original factor, without changing any coordinate of (4).  Only one marginal law is replaced, so the joint law remains a product law and the semi-Markovian independence condition is preserved.

Enumerate the finite set \(\mathcal J\) and perform this replacement successively.  At each step define (2) using the factor laws resulting from the preceding steps.  Equation (4) shows inductively that every step preserves the entire full-data law and every interventional outcome coordinate.  Thus, after the last step,
\[
 P(V=v,R=r)=P_{\mathcal M}(v,r)
 \quad\text{for every }(v,r),
 \qquad
 P(Y=y\mid\operatorname{do}(A=a))=Q_a^{\mathcal M}(y)
 \quad\text{for every }(a,y).
 \tag{5}
\]

It remains to compare the sharp support count with the declared response-function budget.  Let
\[
 S=\sum_{Z\in\mathcal V\cup\mathcal R}|\mathcal X_Z|.
\]
Every alphabet is nonempty, and the distinct values \(a_0,a_1\in\mathcal X_A\) imply \(S\geq2\).  Since \(n\leq2^n\) for every positive integer \(n\), and since both \(|\mathcal X_A|\) and \(|\mathcal X_Y|\) are at most \(S\),
\[
 L\leq2^S,
 \qquad
 K^{\sharp}\leq L+|\mathcal X_A||\mathcal X_Y|
 \leq2^S+S^2.
\]
For integers \(S\geq2\), the elementary inequalities \(S\leq2^{S-1}\) and \(2S-1\leq2^S\) yield
\[
 K^{\sharp}
 \leq2^{2S-2}+2^{2S-2}
 =2^{2S-1}
 \leq2^{2^S}.
 \tag{6}
\]
In the definition of \(N_{\mathcal G}\), every parent-cardinality and missingness-cardinality multiplier is at least one.  Hence \(N_{\mathcal G}\geq S\), and (6) gives
\[
 K^{\sharp}\leq2^{2^S}\leq2^{2^{N_{\mathcal G}}}=K_{\mathcal G}.
 \tag{7}
\]

Relabel the selected support of each factor by \(\{1,\ldots,k_h\}\), where \(k_h\leq K^{\sharp}\), and pad it to \(\{1,\ldots,K_{\mathcal G}\}\) with zero-probability states.  Restrict each original structural equation to the selected factor values.  On any input involving a padded state, extend it by a value constant in all observed parents.  The original restriction respects every label-specific parent omission by the label-semantics assumption, and the constant extension does also; hence the resulting tuple belongs to \(\mathcal T_{\mathcal G}\).  Select that single tuple with \(\gamma_g=1\), set all other one-hot coordinates to zero, and use the padded marginal masses as the \(\rho_{h,k}\).  This is a label-compatible deterministic response representation with mutually independent factors, not a mixture over graph-wide tables.  Its deterministic acyclic evaluation agrees with the reduced model on every positive-probability factor state, so (5) proves both claimed equalities.  No regularity, maximality, positivity, division, or limiting assumption is used.